

\SourceLeaf{221}%
Weil $a$, $b$, $c$ in Bezug auf diesen Halbkreis als verschwindend und die Punkte $A$, $B$ somit als mit $C$ zusammenfallend zu betrachten sind, ist dieser Halbkreis das Maß der drei Winkel des Dreiecks. Zusammengenommen unterscheiden sie sich folglich um weniger als jede gegebene Größe von $180^0$.

Mir scheint, wenn man den Begriff des endlos Wachsenden zulässt, zeigt dieser Beweis sehr einfach, dass in jedem \textls{endlichen} geradlinigen Dreieck die Winkelsumme $=180^0$ ist. Genauer zeigt er, dass die Konstante, die man, wenn Euklids Geometrie nicht wahr wäre, zur Winkelsumme addieren müsste, um die Gleichheit mit $180^0$ herzustellen, kleiner als jede gegebene Größe ist. Da sich dies für jedes Dreieck beweisen lässt, kann diese Konstante ebenso wenig von der Größe des Dreiecks abhängen. \mbox{.\ .\ .\ .\ .\ .}\}

\bigskip
\hrule
\bigskip

\Needspace{7\baselineskip}
\begin{center}
[\textsc{Schumacher} \emph{an} \textsc{Gauss}. \emph{Altona, 29. Juni 1831.}]
\end{center}

\{\mbox{.\ .\ .\ .\ .\ .}\ .\ .\ Nur eines habe ich in Ihrem Brief vermisst, Ihr Urteil über meinen Beweis, dass sich die Winkelsumme eines geradlinigen Dreiecks nur um eine Größe, die kleiner als jede gegebene Größe ist, von $180^0$ unterscheidet. Sie können sich leicht denken, dass mir Ihr Urteil sehr wichtig ist, da Sie jede Schwäche eines Beweises so leicht entdecken. Außer Ihnen, meinem Gehilfen und Professor \textsc{Hansen} vom Seeberg habe ich noch niemandem davon erzählt. Keiner von uns kann einen Paralogismus entdecken.

Sollte jemand einen Beweis dafür verlangen, dass man die Eckpunkte eines endlichen Dreiecks als zusammenfallende Mittelpunkte eines Kreises von unendlichem (brevitatis causa \textls{unendlich} genannt) Radius betrachten kann, so lässt sich dieser Beweis streng führen, auch wenn ich ihn nicht für nötig halte.

Wenn zwei Punkte eine endliche Entfernung voneinander haben, muss man diese Entfernung in Bezug auf eine unendliche Linie, wie mir scheint, $=0$ setzen. In Bezug auf diese unendliche Linie betrachtet, fallen die Punkte somit zusammen. \mbox{.\ .\ .\ .\ .\ .}\ .\}

\bigskip
\hrule
\bigskip
